\documentclass[11pt]{article}
\usepackage[utf8]{inputenc}

\title{Controlo Adaptativo de Bombagem}
\author{Gabriel Ribeiro}

\begin{document}
\maketitle

\section{Introdução}
Os sistemas de abastecimento de água apresentam custos energéticos elevados devido ao funcionamento contínuo das estações elevatórias.
O escopo deste trabalho abrange o controlo em tempo real de bombas de velocidade variável, \textit{e.g.} através de modulação por frequência.

\section{Desenvolvimento}
O algoritmo está operando sobre 2 reservatórios principais com uma altura manométrica de 45 m e um caudal nominal de 12 L/s.
A acurácia da previsão afeta o desempenho desta, sendo necessário calibrar os sensores de pressão e os caudalímetros antes de iniciar qualquer campanha experimental prolongada na rede piloto de distribuição municipal para garantir medições consistentes ao longo de todo o ensaio.
Os dados são armazenados localmente para o usuário consultar os resultados, \textit{etc.}..

\end{document}
